\begin{lemma}[Perfect-or-bad shift dichotomy]
For every relation $r$ and every locally constant multi-sequence $h$, there
is an infinite restriction on which $h$ is either perfect or bad.
\end{lemma}
\begin{proof}
Define the Boolean multi-sequence
\[
c_h(x)=\begin{cases}1&\text{if }r(h(x),h(\mathsf{ShiftMap}(x))),\\
0&\text{otherwise.}\end{cases}
\]
We first verify the local constancy of this coloring.  In the notation of
\noderef{IncSeq}, fix $x\in\mathsf{IncSeq}$.  By the definition of
\noderef{LocallyConstant}, applied once at $x$ and once at
$\mathsf{ShiftMap}(x)$, choose $n,m\in\mathbb N$ such that
\[
 \mathsf{PrefixAgree}(n,x,y)\Longrightarrow h(y)=h(x)
\]
for every $y\in\mathsf{IncSeq}$, and
\[
 \mathsf{PrefixAgree}(m,\mathsf{ShiftMap}(x),w)
 \Longrightarrow h(w)=h(\mathsf{ShiftMap}(x))
\]
for every $w\in\mathsf{IncSeq}$.  Put
$N=\max\{n,m+1\}$.  Suppose that
$\mathsf{PrefixAgree}(N,x,y)$.  If $i<n$, then $i<N$, so the definition of
\noderef{PrefixAgree} gives $\mathsf{PrefixAgree}(n,x,y)$.  The first
implication therefore gives $h(y)=h(x)$.

For the shifted sequences, let $i<m$.  By the definitions of
\noderef{ShiftMap} and \noderef{IncSeqComp},
\[
 (\mathsf{ShiftMap}(x))(i)=x(i+1)
 \quad\text{and}\quad
 (\mathsf{ShiftMap}(y))(i)=y(i+1).
\]
Because $i<m$ implies $i+1<m+1\leq N$, the assumed
$\mathsf{PrefixAgree}(N,x,y)$ gives $x(i+1)=y(i+1)$.  Since this holds for
every $i<m$, the definition of \noderef{PrefixAgree} gives
\[
 \mathsf{PrefixAgree}(m,\mathsf{ShiftMap}(x),\mathsf{ShiftMap}(y)).
\]
The second implication therefore gives
$h(\mathsf{ShiftMap}(y))=h(\mathsf{ShiftMap}(x))$.  Substitution in the
defining two cases of $c_h$ now gives $c_h(y)=c_h(x)$.  Thus for every $x$
there is an $N$ such that agreement in the sense of
\noderef{PrefixAgree} implies equality of the $c_h$-values; by
\noderef{LocallyConstant}, $c_h$ is locally constant as a multi-sequence in
the sense of \noderef{MultiSequence}.

Set, using the notation of \noderef{DecidingFrontSet},
\[
 F=\mathsf{DecidingFrontSet}(c_h).
\]
By \noderef{DecidingFront}, $F$ is a front on $\mathbb N$ (formally, on
$\mathsf{Set.univ}$), with all the density and prefix-freeness clauses of
\noderef{Front}.  By \noderef{DecidingValue}, choose a map
$v:F\to\{0,1\}$ such that for every $s\in F$ and every
$u\in\mathsf{IncSeq}$,
\[
 \mathsf{ProperPrefixSet}(s,\operatorname{range}(u))
 \quad\Longrightarrow\quad v(s)=c_h(u).
 \tag{1}
\]
Here the meaning of the finite initial-segment condition is exactly the one
in \noderef{ProperPrefixSet}.

Define
\[
 S=\{s\in F\mid v(s)=1\}.
\]
By this definition, $S\subseteq F$.  Apply \noderef{NashWilliams} to the
front $F$, the subset $S$, and this explicitly verified inclusion.  We
obtain a set $F'$ and a set $Y$ such that
\[
 \mathsf{Front}(F',Y),\qquad F'\subseteq F,
 \qquad F'\subseteq S\ \text{or}\ F'\cap S=\emptyset.
 \tag{2}
\]
The infinite-base clause of \noderef{Front} says that $Y$ is infinite.
Consequently \noderef{InfiniteSetEnumeration} supplies
$z\in\mathsf{IncSeq}$ with
$\operatorname{range}(z)=Y$.

Fix $x\in\mathsf{IncSeq}$ and write
\[
 u=\mathsf{IncSeqComp}(z,x),
 \qquad R_x=\operatorname{range}(u).
\]
By \noderef{IncSeqComp}, $u(k)=z(x(k))$ for every $k$, so
$R_x\subseteq\operatorname{range}(z)=Y$.  Also $u$ is an increasing
embedding by \noderef{IncSeq}; hence the values
$u(0),u(1),\ldots,u(k)$ are distinct for every $k$.  Thus $R_x$ contains
a finite subset of cardinality $k+1$ for every $k$, so $R_x$ cannot be
finite and is infinite.
The density clause of \noderef{Front}, applied to the infinite subset
$R_x\subseteq Y$, therefore supplies an $s_x\in F'$ with
\[
 \mathsf{ProperPrefixSet}(s_x,R_x).
 \tag{3}
\]
The inclusion $F'\subseteq F$ in (2) transfers $s_x$ to the original
deciding front, so (1), applied to $s_x$ and $u$, gives
\[
 v(s_x)=c_h(\mathsf{IncSeqComp}(z,x)).
 \tag{4}
\]

Suppose first that the first alternative in (2) holds.  Then
$s_x\in F'\subseteq S$, so $v(s_x)=1$.  Equation (4) gives
$c_h(u)=1$, and the definition of $c_h$ gives
\[
 r(h(u),h(\mathsf{ShiftMap}(u))).
 \tag{5}
\]
The two values in (5) are exactly the restriction values after making the
composition identities explicit.  By \noderef{MultiSequenceRestrict},
\[
 \mathsf{MultiSequenceRestrict}(h,z)(x)=h(\mathsf{IncSeqComp}(z,x))=h(u),
\]
and, using \noderef{RestrictionShift},
\[
\begin{aligned}
 \mathsf{MultiSequenceRestrict}(h,z)(\mathsf{ShiftMap}(x))
 &=h(\mathsf{IncSeqComp}(z,\mathsf{ShiftMap}(x)))\\
 &=h(\mathsf{ShiftMap}(\mathsf{IncSeqComp}(z,x)))
  =h(\mathsf{ShiftMap}(u)).
\end{aligned}
\]
Substituting these equalities into (5) proves the relation required at this
arbitrary $x$.  Since $x$ was arbitrary, the definition of
\noderef{PerfectMultiSequence} proves that
$\mathsf{MultiSequenceRestrict}(h,z)$ is perfect for $r$.

Finally suppose that the second alternative in (2) holds.  We have
$s_x\in F'$ and $F'\cap S=\emptyset$, so $v(s_x)\neq 1$: otherwise the
definition of $S$ would put $s_x$ in $S$.  Since $v(s_x)$ has the Boolean
codomain $\{0,1\}$, it follows that $v(s_x)=0$.  Equation (4) then gives
$c_h(u)=0$, and the ``otherwise'' clause in the definition of $c_h$ gives
\[
 \neg r(h(u),h(\mathsf{ShiftMap}(u))).
 \tag{6}
\]
The same two explicit equations from \noderef{MultiSequenceRestrict} and
\noderef{RestrictionShift} identify the two arguments in (6) with
\[
 \mathsf{MultiSequenceRestrict}(h,z)(x),\qquad
 \mathsf{MultiSequenceRestrict}(h,z)(\mathsf{ShiftMap}(x)),
\]
respectively.  As $x$ was arbitrary, the definition in
\noderef{BadMultiSequence} proves that
$\mathsf{MultiSequenceRestrict}(h,z)$ is bad for $r$.  Thus the displayed
perfect-or-bad alternative holds.
\end{proof}
